\documentclass[12pt]{article}
\usepackage{fontspec}
\setmainfont{EB Garamond}[Numbers=OldStyle, Ligatures=TeX]
\usepackage[letterpaper, margin=1.25in]{geometry}
\usepackage{hyperref}
\usepackage{orcidlink}

\begin{document}
\thispagestyle{empty}

\begin{center}
{\Large\bfseries What Dogs Know About Trails:\\
Metaphor, Mechanisms, and the Phylogeny of Cross-Domain Extension}

\bigskip\bigskip

Brett Reynolds \orcidlink{0000-0003-0073-7195}

\medskip

Humber Polytechnic \& University of Toronto

\medskip

\href{mailto:brett.reynolds@humber.ca}{brett.reynolds@humber.ca}
\end{center}

\bigskip

\noindent\textbf{Keywords:} conceptual metaphor, homeostatic property cluster, cross-domain extension, categorization, canine cognition, theory-constitutive metaphor, projectibility

\bigskip

\noindent\textbf{Corresponding author:} Brett Reynolds, Humber Polytechnic, 205 Humber College Blvd, Toronto, ON M9W 5L7, Canada. Email: \href{mailto:brett.reynolds@humber.ca}{brett.reynolds@humber.ca}

\bigskip

\noindent\textbf{Acknowledgements:} I drafted this paper with the assistance of Claude Code (Opus 4.6) and have reviewed and revised all content. I take full responsibility for the final text.

\bigskip

\noindent\textbf{Conflict of interest:} None.

\bigskip

\noindent\textbf{Data availability:} This is a theoretical paper with no empirical data. The preprint is available at \url{https://osf.io/preprints/psyarxiv/vuafy}.

\end{document}
